% contents/chapter-03/section-09.tex 第三章第三节
\documentclass[../../main.tex]{subfiles}% 注意这里的文件路径不能用 ./main.tex ,否则用latexmk编译子文件会报错
\graphicspath{{\subfix{./image/}}{\subfix{../../image/}}} % 指定图片引用路径,后续可以直接使用图片文件名
% 如果图片存放在上级目录,则使用 ../../image/ 作为备用路径

% 例如:
% \begin{figure}[H]
% \centering
% \includegraphics[width=0.3\textwidth]{图.png}
% \caption{}
% \label{figure:图}
% \end{figure}
% 注意:上述\label{}一定要放在\caption{}之后,否则引用图片序号会只会显示??.

\begin{document}

\section{覆盖定理}

\begin{proposition}\label{proposition:真子空间至多包含n-1个基向量}
设\(V_0\)是数域\(\mathbb{F}\)上$n$维向量空间\(V\)的真子空间,则\(V_0\)至多包含$n-1$个$V$中的基向量.
\end{proposition}
\begin{proof}
反证法,若$V_0$包含\(n\)个\(V\)中的基向量,则$V_i$就包含了$V$的一组基.不妨设$V_0$中的这组基向量为$\{e_1,e_2,\cdots e_n\}$,则$\forall \alpha\in V$,有$\alpha =k_1e_1+k_2e_2+\cdots +k_ne_n \in V_0$,其中$k_i\in \mathbb{F}$,$i=1,2,\cdots,n$.故$V_0\supset V$,又$V_0\subset V$,因此$V_0=V$.这与\(V_0\)是\(V\)的真子空间矛盾.
\end{proof}

\begin{theorem}[覆盖定理]\label{theorem:覆盖定理}
\begin{enumerate}[(1)]
\item\label{theorem:覆盖定理-1} 数域 \( \mathbb{F} \)上 \( n \) 维线性空间 \( V \) 不能分解为有限个 (大于等于 2 个) 非平凡子空间之并.  

\item\label{theorem:覆盖定理-2} 域 \( \mathbb{F} \) 上线性空间 \( V \) 不能分解为两个非平凡子空间之并. 

\item\label{theorem:覆盖定理-3} 域 \( \mathbb{F} \) 上无限维线性空间 \( V \) 不能分解为有限个 (大于等于 2 个) 非平凡子空间之并.  
\end{enumerate}
\end{theorem}
\begin{remark}
这个定理表明:\textbf{有限个真子空间不能覆盖全空间}.
下述证明要用到任意一个数域都有无穷个元素这一事实. 因此,对于有限域上的向量空间,上例结论不一定成立.
\end{remark}
\begin{proof}
\begin{enumerate}[(1)]
\item {\color{blue}证法一:}设\(V_1,V_2,\cdots,V_m\)是数域\(\mathbb{F}\)上向量空间\(V\)的\(m\)个真子空间,则只需证明在\(V\)中必存在一个向量\(\boldsymbol{\alpha}\),它不属于任何一个\(V_i\).
对个数\(m\)进行归纳,当\(m = 1\)时结论显然成立. 设\(m = k\)时结论成立,现要证明\(m = k + 1\)时结论也成立. 由归纳假设,存在向量\(\boldsymbol{\alpha}\),它不属于任何一个\(V_i(1\leqslant  i\leqslant  k)\). 若\(\boldsymbol{\alpha}\)也不属于\(V_{k + 1}\),则结论已成立,因此可设\(\boldsymbol{\alpha}\in V_{k + 1}\). 在\(V_{k + 1}\)外选一个向量\(\boldsymbol{\beta}\),作集合
\begin{align*}
M = \{t\boldsymbol{\alpha}+\boldsymbol{\beta}|t\in\mathbb{F}\}.
\end{align*}
事实上,我们可将\(M\)看成是通过\(\boldsymbol{\beta}\)的终点且平行于\(\boldsymbol{\alpha}\)的一根“直线”,现要证明它和每个\(V_i\)最多只有一个交点. 首先,\(M\)和\(V_{k + 1}\)无交点,因为若\(t\boldsymbol{\alpha}+\boldsymbol{\beta}\in V_{k + 1}\),则从\(t\boldsymbol{\alpha}\in V_{k + 1}\)可推出\(\boldsymbol{\beta}\in V_{k + 1}\),与假设矛盾. 又若对某个\(V_i(i<k + 1)\),存在\(t_1\neq t_2\),使得\(t_1\boldsymbol{\alpha}+\boldsymbol{\beta}\in V_i\),\(t_2\boldsymbol{\alpha}+\boldsymbol{\beta}\in V_i\),则\((t_1 - t_2)\boldsymbol{\alpha}\in V_i\),从而导致\(\boldsymbol{\alpha}\in V_i\),与假设矛盾. 因此,\(M\)和每个\(V_i\)最多只有一个交点,从而\(M\)中只有有限个向量属于\(\bigcup_{i=1}^{k+1}{V_i}\),而\(t\)有无穷多个选择,即$M$中有无穷多个向量.故必存在$\gamma\in M\setminus \bigcup_{i=1}^{k+1}{V_i}$,由此即得结论. 

{\color{blue}证法二:}设\(V_1,V_2,\cdots,V_m\)是数域\(\mathbb{F}\)上向量空间\(V\)的\(m\)个真子空间,则只需证明在\(V\)中必存在一个向量\(\boldsymbol{\alpha}\),它不属于任何一个\(V_i\).任取\(V\)的一组基\(\{\boldsymbol{e}_1,\boldsymbol{e}_2,\cdots,\boldsymbol{e}_n\}\). 对任意的正整数\(k\),构造\(V\)中向量
\begin{align*}
\boldsymbol{\alpha}_k=\boldsymbol{e}_1 + k\boldsymbol{e}_2+\cdots + k^{n - 1}\boldsymbol{e}_n.
\end{align*}
设向量族\(S = \{\boldsymbol{\alpha}_k|k = 1,2,\cdots\}\). 由\hyperref[example:3.110.1]{例题\ref{example:3.110.1}}可知,\(S\)中任意\(n\)个不同的向量都构成\(V\)的一组基. 因为\(V_i\)都是\(V\)的真子空间,所以每个\(V_i\)至多包含\(S\)中\(n - 1\)个向量.因此$\bigcup_{i=1}^m{V_i}$至多包含$S$中$m(n-1)$个向量.
又由于\(S\)是无限集合,故存在某个向量\(\boldsymbol{\alpha}_k\),使得\(\boldsymbol{\alpha}_k\)不属于任何一个\(V_i\).

{\color{blue}证法三:}设$V=\bigcup_{i=1}^m{V_i}$,其中$m\geqslant 2$,$V_i$都是非平凡子空间.考虑 \( V \) 的基 \( e_1, e_2, \cdots, e_n \) 和一组向量  
\[
\alpha_i = e_1 + i e_2 + i^2 e_3 + \cdots + i^{n-1} e_n, \quad i = 1, 2, \cdots.
\] 
从而
\begin{align*}
\left( \alpha _{i_1},\alpha _{i_2},\cdots ,\alpha _{i_n} \right) =\left( e_1,e_2,\cdots ,e_n \right) \left( \begin{matrix}
1&		1&		\cdots&		1\\
i_1&		i_2&		\cdots&		i_n\\
\vdots&		\vdots&		&		\vdots\\
i_{1}^{n-1}&		i_{2}^{n-1}&		\cdots&		i_{n}^{n-1}\\
\end{matrix} \right) ,i_1,\cdots ,i_n\in \mathbb{N} .
\end{align*}
由Vandermonde行列式的性质, 上面这组向量$\{\alpha_i\}_{i=1}^{\infty}$中任何 \( n \) 个都线性无关. 而其中必有无限个$\alpha_i$落入某个题设 \( V \) 的某个非平凡子空间$V_i$中,否则,$V=\bigcup_{i=1}^nV_i$中只包含$\{\alpha_i\}_{i=1}^{\infty}$中有限个向量,这与$\{\alpha_i\}_{i=1}^{\infty}\subset V$矛盾! 从而这个子空间至少是\( n \)维的, 这和其是非平凡子空间矛盾!  

\item 设 \( V = V_1 \bigcup V_2 \), \( V_1, V_2 \) 都是 \( V \) 的非平凡不变子空间. 则存在 \( v_1 \in V_1 \setminus V_2, v_2 \in V_2 \setminus V_1 \). 注意到
\[
v_1 + v_2 \in V_1 \Rightarrow v_2 \in V_1, \, v_1 + v_2 \in V_2 \Rightarrow v_1 \in V_2,
\]  
从而  
\[
v_1 + v_2 \notin V_1 \bigcup V_2,
\]  
这就是一个矛盾! 因此我们证明了域 \( \mathbb{F} \) 上线性空间 \( V \) 不能分解为两个非平凡子空间之并.  

\item 若 \( V_i \subset V, i = 1, 2, \cdots, n \) 是非平凡子空间. 我们归纳证明  
\begin{align}
V \neq \bigcup_{i=1}^n V_i. \label{eq:::::---64486-23.21}
\end{align}
当 \( n = 1 \) 显然有 \(\eqref{eq:::::---64486-23.21}\) 成立. 假设对 \( n \) 有 \(\eqref{eq:::::---64486-23.21}\) 成立, 当 \( n + 1 \) 时, 由归纳假设取 \( \alpha \in V \setminus \left( \bigcup_{i=1}^n V_i \right) \).  

如果 \( \alpha \notin V_{n+1} \), 则 \(\eqref{eq:::::---64486-23.21}\) 对 \( n + 1 \) 已经成立.  
- 若 \( \alpha \in V_{n+1} \), 由 \( V_{n+1} \) 非平凡, 存在 \( \beta \in V \setminus V_{n+1} \). 注意到  
\[
S = \{ t\alpha + \beta : t \in \mathbb{F} \} \subset V \setminus V_{n+1},
\]  
这是因为若某个 \( t\alpha + \beta \in V_{n+1} \) 可推出 \( \beta \in V_{n+1} \) 而矛盾!. 现在若有 \( t_1\alpha + \beta, t_2\alpha + \beta \) 属于同一个 \( V_j, j \in \{1, 2, \cdots, n\} \), 则  
\[
t_1\alpha + \beta - (t_2\alpha + \beta) = (t_1 - t_2) \alpha \in V_j,
\]  
即 \( t_1 = t_2 \). 现在 \( S \) 中向量有无穷多个, 所以必然有一个不落在 \( \bigcup_{j=1}^n V_j \), 自然也不落在 \( \bigcup_{j=1}^{n+1} V_j \), 这就证明了 \(\eqref{eq:::::---64486-23.21}\) 对 \( n + 1 \) 也成立!.  
至此我们完成了证明.
\end{enumerate}
\end{proof}

\begin{corollary}\label{corollary:真子空间外仍有一组基存在}
设\(V_1,V_2,\cdots,V_m\)是数域\(\mathbb{F}\)上向量空间\(V\)的\(m\)个真子空间,证明:\(V\)中必有一组基,使得每个基向量都不在诸\(V_i\)的并中.
\end{corollary}
\begin{proof}

{\color{blue}证法一:}
由\hyperref[theorem:覆盖定理]{覆盖定理}可知,存在非零向量\(\boldsymbol{e}_1\in V\),使得\(\boldsymbol{e}_1\notin\bigcup_{i = 1}^{m}V_i\). 定义\(V_{m + 1}=L(\boldsymbol{e}_1)\),再由\hyperref[theorem:覆盖定理]{覆盖定理}可知,存在向量\(\boldsymbol{e}_2\in V\),使得\(\boldsymbol{e}_2\notin\bigcup_{i = 1}^{m + 1}V_i\). 由\hyperref[corollary:线性无关向量组的命题1]{推论\ref{corollary:线性无关向量组的命题1}}可知,\(\boldsymbol{e}_2\notin L(\boldsymbol{e}_1)\)意味着\(\boldsymbol{e}_1,\boldsymbol{e}_2\)线性无关. 重新定义\(V_{m + 1}=L(\boldsymbol{e}_1,\boldsymbol{e}_2)\),再由\hyperref[theorem:覆盖定理]{覆盖定理}可知,存在向量\(\boldsymbol{e}_3\in V\),使得\(\boldsymbol{e}_3\notin\bigcup_{i = 1}^{m + 1}V_i\). 再由\hyperref[corollary:线性无关向量组的命题1]{推论\ref{corollary:线性无关向量组的命题1}}可知,\(\boldsymbol{e}_3\notin L(\boldsymbol{e}_1,\boldsymbol{e}_2)\)意味着\(\boldsymbol{e}_1,\boldsymbol{e}_2,\boldsymbol{e}_3\)线性无关. 不断重复上述讨论,即添加线性无关的向量重新定义\(V_{m + 1}\),并反复利用\hyperref[theorem:覆盖定理]{覆盖定理}和\hyperref[corollary:线性无关向量组的命题1]{推论\ref{corollary:线性无关向量组的命题1}}的结论,最后可以得到\(n\)个线性无关的向量\(\boldsymbol{e}_1,\boldsymbol{e}_2,\cdots,\boldsymbol{e}_n\),它们构成\(V\)的一组基,且满足\(\boldsymbol{e}_j\notin\bigcup_{i = 1}^{m}V_i(1\leqslant  j\leqslant  n)\). 

{\color{blue}证法二:}任取\(V\)的一组基\(\{\boldsymbol{e}_1,\boldsymbol{e}_2,\cdots,\boldsymbol{e}_n\}\). 对任意的正整数\(k\),构造\(V\)中向量\(\boldsymbol{\alpha}_k=\boldsymbol{e}_1 + k\boldsymbol{e}_2+\cdots + k^{n - 1}\boldsymbol{e}_n\),设向量族\(S = \{\boldsymbol{\alpha}_k|k = 1,2,\cdots\}\). 由\hyperref[example:3.110.1]{例题\ref{example:3.110.1}}可知,\(S\)中任意\(n\)个不同的向量都构成\(V\)的一组基. 因为\(V_i\)都是\(V\)的真子空间,所以每个\(V_i\)至多包含\(S\)中\(n - 1\)个向量.因此$\bigcup_{i=1}^m{V_i}$至多包含$S$中$m(n-1)$个向量.
又由于\(S\)是无限集合,故存在某个向量\(\boldsymbol{\alpha}_k\),使得\(\boldsymbol{\alpha}_k\)不属于任何一个\(V_i\).
进一步,在\(S\)中一定还存在\(n\)个不同的向量\(\boldsymbol{\alpha}_{k_1},\boldsymbol{\alpha}_{k_2},\cdots,\boldsymbol{\alpha}_{k_n}\),使得每个\(\boldsymbol{\alpha}_{k_j}\)都不属于任何一个\(V_i\),此时\(\{\boldsymbol{\alpha}_{k_1},\boldsymbol{\alpha}_{k_2},\cdots,\boldsymbol{\alpha}_{k_n}\}\)就构成了\(V\)的一组基.
\end{proof}

\begin{example}
设 \(V\) 是复数域上线性空间且 \(V_i, i=1,2,\cdots,k\) 是 \(V\) 真子空间. 若 \(\alpha_i \in V, i=1,2,\cdots,k\), 证明: \(V \neq \bigcup\limits_{i=1}^k (\alpha_i + V_i).\)
\end{example}
\begin{proof}

{\color{blue}证法一:}选取特殊向量 \(v_1\): 因为 \(V_i\) 都是 \(V\) 的真子空间, 根据\hyperref[theorem:覆盖定理]{覆盖定理}，我们知道存在向量 \(v_1 \in V\)，使得 \(v_1 \notin V_i\) 对所有的 \(i=1,2,\cdots,k\) 都成立.
任取一个固定的向量 \(v_0 \in V\)，构造直线
\begin{align*}
L = \{ v_0 + t v_1 \mid t \in \mathbb{C} \}.
\end{align*}
因为 \(\mathbb{C}\) 是无限域，所以直线 \(L\) 上有无穷多个点.
下证直线 \(L\) 与 \(\alpha_i + V_i\) 的交点至多只有一个.假设存在\(t_1 \neq t_2\)，使得
\begin{align*}
v_0 + t_1 v_1,v_0 + t_2 v_1 \in \alpha_i + V_i.
\end{align*}
从而$(t_1 - t_2)v_1 \in V_i.$
故\(v_1 \in V_i\). 这与\(v_1\)的取法矛盾！
因此，对于每一个 \(i\)，直线 \(L\) 与 \(\alpha_i + V_i\) 的交点最多只有一个.又因为直线 \(L\) 包含了无穷多个点,所以一定存在$t_0\in \mathbb{C}$,使得
\begin{align*}
v_0 + t_0 v_1 \in L\setminus \bigcup\limits_{i=1}^k (\alpha_i + V_i)\subset V\setminus \bigcup\limits_{i=1}^k (\alpha_i + V_i). 
\end{align*}
因此
\begin{align*}
V \neq \bigcup\limits_{i=1}^k (\alpha_i + V_i).
\end{align*}

{\color{blue}证法二:}考虑线性空间
\begin{align*}
\tilde{V} \triangleq V \oplus \mathbb{C}, \quad \tilde{V}_i \triangleq \{ (\beta + t\alpha_i, t) : \beta \in V_i, t \in \mathbb{C} \}, \quad i=1,2,\cdots,k.
\end{align*}
对 \(i=1,2,\cdots,k\), 注意到对 \(\alpha \in V \setminus V_i\), 我们有 \((\alpha, 0) \notin \tilde{V}_i\), 即 \(\tilde{V}_i\) 都是 \(\tilde{V}\) 的真子空间. 由\hyperref[theorem:覆盖定理]{覆盖定理}得
\begin{align*}
\tilde{V} \neq (V \times \{0\}) \cup \left( \bigcup\limits_{i=1}^k \tilde{V}_i \right).
\end{align*}
故存在
\begin{align*}
(\bm{x}, t) \in \tilde{V} \setminus \left( (V \times \{0\}) \cup \left( \bigcup\limits_{i=1}^k \tilde{V}_i \right) \right), 
\end{align*}
从而
\begin{align*}
t \neq 0, \quad (\bm{x}, t) = \left( t\left(\frac{\bm{x}}{t} - \alpha_i\right) + t\alpha_i, t \right) \notin \tilde{V}_i, \quad i=1,2,\cdots,k.
\end{align*}
现在就有
\begin{align*}
\frac{\bm{x}}{t} - \alpha_i \notin V_i, \quad i=1,2,\cdots,k, 
\end{align*}
即
\begin{align*}
\frac{\bm{x}}{t} \in V\setminus \bigcup\limits_{i=1}^k (\alpha_i + V_i). 
\end{align*}
因此
\begin{align*}
V \neq \bigcup\limits_{i=1}^k (\alpha_i + V_i).
\end{align*}
\end{proof}

\begin{example}
设 \( V \) 是有限维线性空间，\( A_1, A_2, \cdots, A_s \) 是 \( V \) 上两两不同线性变换，则 \( \exists \alpha \in V \)，使 \( A_1\alpha, A_2\alpha, \cdots, A_s\alpha \) 两两不同.
\end{example}
\begin{proof}
考虑$\bigcup_{1\leqslant i<j\leqslant n}{\mathrm{Ker}\left( A_j-A_i \right)}$，显然$\mathrm{Ker}\left( A_j-A_i \right) \ne V,\varnothing$。故由\hyperref[theorem:覆盖定理]{覆盖定理}知，我们有$\bigcup_{1\leqslant i<j\leqslant n}{\mathrm{Ker}\left( A_j-A_i \right)}\ne V$。因此$\exists \alpha \in V\backslash \bigcup_{1\leqslant i<j\leqslant n}{\mathrm{Ker}\left( A_j-A_i \right)}$，此时
\begin{align*}
\left( A_j-A_i \right) \alpha \ne 0\Longleftrightarrow A_i\alpha \ne A_j\alpha ,1\leqslant i<j\leqslant n.
\end{align*}
\end{proof}

\begin{example}
设 \( T \) 是 \( \mathbb{R}^{n \times n} \) 上的线性变换且 \( T(AB) = T(A)T(B) \)，\( T(AB) = T(B)T(A) \) 必有一个成立. 证明: 要么对所有 \( A, B \in \mathbb{R}^{n \times n} \) 都有 \( T(AB) = T(A)T(B) \)，要么对所有 \( A, B \in \mathbb{R}^{n \times n} \) 都有 \( T(AB) = T(B)T(A) \).  
\end{example}
\begin{proof}
对$\forall A\in \mathbb{R}^{n\times n},$考虑  
\[
V_A = \left\{ B \in \mathbb{R}^{n \times n} : T(AB) = T(B)T(A) \right\};
\]  
\[
V_A' = \left\{ B \in \mathbb{R}^{n \times n} : T(AB) = T(A)T(B) \right\},
\]  
则有 \( \mathbb{R}^{n \times n} = V_A \bigcup V_A' \). 由\hyperref[theorem:覆盖定理]{覆盖定理}，我们有 \( \mathbb{R}^{n \times n} = V_A \) 或者 \( \mathbb{R}^{n \times n} = V_A' \).  
再由$A$的任意性可知结论成立.
\end{proof}

\begin{definition}
设 \( V \) 是域 \( \mathbb{F} \) 上的有限维线性空间，\( v \in V \)，\( A \) 是 \( V \) 上线性变换。记 \( m_v \in \mathbb{F}[x] \) 是 \( m_v(A)v = 0 \) 成立的最低次的首一多项式，我们叫 \( m_v \) 为\textbf{ \( \boldsymbol{v} \) 的极小多项式}。
\end{definition}

\begin{proposition}\label{proposition:向量的极小多项式的整性}
设 \( V \) 是域 \( \mathbb{F} \) 上的有限维线性空间，\( v \in V \)，\( A \) 是 \( V \) 上线性变换，\( m_v \) 为 \( v \) 的极小多项式。若$m\in \mathbb{F}[x]$,且$m(A)v=0$,则必有$m_v\mid m.$
\end{proposition}
\begin{proof}
由带余除法可知，存在\( q,r\in \mathbb{F}[x] \)，且\( \deg r<\deg m_v \)，使得\( m=qm_v+r \)。从而
\begin{align*}
0=m(A)v=q(A)m_v(A)v+r(A)v=r(A)v.
\end{align*}
于是\( r\equiv 0 \)，否则与\( m_v \)是\( v \)的极小多项式矛盾！故\( m_v|m \)。
\end{proof}

\begin{example}
设 \( V \) 是域 \( \mathbb{F} \) 上的有限维线性空间，\( v \in V \)，\( A \) 是 \( V \) 上线性变换。记 \( m_v \in \mathbb{F}[x] \) 是 \( m_v(A)v = 0 \) 成立的最低次的首一多项式，我们叫 \( m_v \) 为 \( v \) 的极小多项式。证明：存在 \( v \in V \) 使得 \( v \) 的极小多项式恰好是 \( A \) 的极小多项式。
\end{example}
\begin{proof}
设 \( m \in \mathbb{F}[x] \) 是 \( A \) 的极小多项式，那么
\[
m(A)v = 0, \forall v \in V
\]
{\color{blue}证法一:}
反证,假设对$\forall v\in V$,都有$m_v(x)\ne m(x)$.
又由\refpro{proposition:向量的极小多项式的整性}知\( m_v|m \)。因此$m_v(x)(v\in V)$都是$m(x)$的首一真因式.
考虑 \( m \) 的所有首一真因式
\[
m_1, m_2, \cdots, m_N \in \mathbb{F}[x],N\in\mathbb{N},
\]
则对$\forall v\in V$,都有$m_v\in \{m_1,m_2,\cdots,m_N\}$.
考虑
\[
V_i \triangleq \{ v \in V : m_i(A)v = 0 \}, i = 1, 2, \cdots, N.
\]
对$\forall v_0\in V,$由于$m_{v_0}(x)$都是$m(x)$的首一真因式,故存在$i_0\in [1,N]\cap \mathbb{N}$,使得$m_{v_0}=m_{i_0}$.从而$m_{i_0}(A)v_0=0$,于是$v_0\in V_{i_0}.$因此$V\subset \bigcup_{i=1}^N V_i.$
故\( V = \bigcup_{i=1}^N V_i \)，于是由\hyperref[theorem:覆盖定理]{覆盖定理}，我们有某个 \( i \) 使得 \( V_i = V \)，故$$m_i(A)v=0,\forall v\in V.$$
因此$m_i(A)=0.$
这意味着 \( m | m_i \)，故 \( m = m_i\)。这与$m_i$是$m$的真因式矛盾!

{\color{blue}证法二:}设$m(x)=(x-\lambda_1)^{k_1}(x-\lambda_2)^{k_2}\cdots(x-\lambda_t)^{k_t}$，则由\refthe{theorem:线性变换互素分解对应核子空间直和分解}知
\begin{align}
\mathbb{F}^n=\operatorname{Ker}(A-\lambda_1I)^{k_1}\oplus\cdots\oplus\operatorname{Ker}(A-\lambda_tI)^{k_t}.\label{eq::--2r89fujfis9fj23ojcoszjfojpj3fgsdhdsw211}
\end{align}
对$\forall i\in\{1,2,\cdots,t\}$，若$\operatorname{Ker}(A-\lambda_iI)^{k_i}=\operatorname{Ker}(A-\lambda_iI)^{k_i-1}$，则
\begin{align*}
\mathbb{F}^n=\operatorname{Ker}(A-\lambda_1I)^{k_1}\oplus\cdots\oplus\operatorname{Ker}(A-\lambda_iI)^{k_i-1}\oplus\cdots\oplus\operatorname{Ker}(A-\lambda_tI)^{k_t}.
\end{align*}
从而对$\forall \alpha\in\mathbb{F}^n$，都存在$\beta_j\in\operatorname{Ker}(A-\lambda_jI)^{k_j}$，$j\in\{1,2,\cdots,t\}\setminus\{i\}$，$\beta_i\in\operatorname{Ker}(A-\lambda_iI)^{k_i-1}$，使得
\begin{align*}
\alpha=\beta_1+\cdots+\beta_i+\cdots+\beta_t.
\end{align*}
令$g(x)=(x-\lambda_1)^{k_1}\cdots(x-\lambda_i)^{k_i-1}\cdots(x-\lambda_t)^{k_t}$，则$\deg g(x)<\deg m(x)$且
\begin{align*}
g(A)\alpha=g(A)\beta_1+\cdots+g(A)\beta_i+\cdots+g(A)\beta_t=0.
\end{align*}
由$\alpha$的任意性知$g(x)$是$A$的零化多项式，因此$m(x)\mid g(x)$，但$\deg g(x)<\deg m(x)$，矛盾！故
\begin{align*}
\operatorname{Ker}(A-\lambda_iI)^{k_i}\supset\operatorname{Ker}(A-\lambda_iI)^{k_i-1}\Longrightarrow\operatorname{Ker}(A-\lambda_iI)^{k_i}\setminus\operatorname{Ker}(A-\lambda_iI)^{k_i-1}\neq\varnothing.
\end{align*}
现在任取$\alpha_i\in\operatorname{Ker}(A-\lambda_iI)^{k_i}\setminus\operatorname{Ker}(A-\lambda_iI)^{k_i-1}$，记$m_i(x)=(x-\lambda_i)^{k_i}$，则
\begin{align*}
m_i(A)\alpha_i=0,\quad i=1,2,\cdots,t.
\end{align*}
故由\refpro{proposition:向量的极小多项式的整性}知，$m_{\alpha_i}(x)\mid m_i(x)$，再由首一性可设$m_{\alpha_i}(x)=(x-\lambda_i)^{m_i}$，$m_i\leqslant k_i$。若$m_i<k_i$，则
\begin{align*}
0=m_{\alpha_i}(A)\alpha_i\Longrightarrow\alpha_i\in\operatorname{Ker}(A-\lambda_iI)^{m_i}\subseteq\operatorname{Ker}(A-\lambda_iI)^{k_i-1},
\end{align*}
这与$\alpha_i$的取法矛盾！故$m_i=k_i$，即$m_{\alpha_i}(x)=(x-\lambda_i)^{k_i}$。进而
\begin{align*}
\left[m_{\alpha_1}(x),m_{\alpha_2}(x),\cdots,m_{\alpha_t}(x)\right]=m(x).
\end{align*}
令$\alpha=\alpha_1+\alpha_2+\cdots+\alpha_t$，则由\refpro{proposition:向量的极小多项式的整性}知$m_{\alpha}(x)\mid m(x)$，并且
\begin{align*}
0=m_{\alpha}(A)\alpha=m_{\alpha}(A)\alpha_1+m_{\alpha}(A)\alpha_2+\cdots+m_{\alpha}(A)\alpha_t.
\end{align*}
又因为$\operatorname{Ker}(A-\lambda_iI)^{k_i}$都是$A$的不变子空间，所以$m_{\alpha}(A)\alpha_i\in\operatorname{Ker}(A-\lambda_iI)^{k_i}$。由\eqref{eq::--2r89fujfis9fj23ojcoszjfojpj3fgsdhdsw211}式和直和的充要条件知
\begin{align*}
m_{\alpha}(A)\alpha_i=0,\quad i=1,2,\cdots,t.
\end{align*}
由\refpro{proposition:向量的极小多项式的整性}知$m_{\alpha_i}(x)\mid m_{\alpha}(x)$，$i=1,2,\cdots,t$。进而$\left[m_{\alpha_1}(x),m_{\alpha_2}(x),\cdots,m_{\alpha_t}(x)\right]\mid m_{\alpha}(x)$，即$m(x)\mid m_{\alpha}(x)$。因此$m(x)=m_{\alpha}(x)$.
\end{proof}

\begin{proposition}\label{proposition:无穷多个线性子空间的交等于有限个的交}
设 \( V \) 是域 \( \mathbb{F} \) 上的有限维线性空间，\(\{V_{\alpha}\}_{\alpha \in \Lambda}\) 是 \( V \) 的子空间. 证明存在 \(\alpha_1, \alpha_2, \cdots, \alpha_n \in \Lambda\) 使得
\[
\bigcap_{i = 1}^n V_{\alpha_i} = \bigcap_{\alpha \in \Lambda} V_{\alpha}
\]
\end{proposition}
\begin{proof}
若 \(|\Lambda| < \infty\)，则已经得证明. 若 \(|\Lambda| = \infty\)，取 \(V_{\alpha_1}\)，若
\[
\dim V_{\alpha_1} = \dim \bigcap_{\alpha \in \Lambda} V_{\alpha}
\]
这已经得证.

若
\[
\dim V_{\alpha_1} > \dim \bigcap_{\alpha \in \Lambda} V_{\alpha}
\]
先证存在 \(\alpha_2 \in \Lambda \setminus \{\alpha_1\}\) ,使得
\[
V_{\alpha_1} \bigcap V_{\alpha_2} \subset V_{\alpha_1}.
\]
假设对任何 \(\alpha \in \Lambda \setminus \{\alpha_1\}\)，我们有
\[
V_{\alpha_1} \bigcap V_{\alpha} = V_{\alpha_1}\Longrightarrow \bigcap_{\alpha \in \Lambda}{V_{\alpha}}=V_{\alpha _1}\cap \left( \bigcap_{\alpha \in \Lambda \backslash \left\{ \alpha _1 \right\}}{V_{\alpha}} \right) =V_{\alpha _1}.
\]
则
\[
\dim V_{\alpha_1} = \dim \bigcap_{\alpha \in \Lambda} V_{\alpha}
\]
矛盾！故存在 \(\alpha_2 \in \Lambda \setminus \{\alpha_1\}\) ,使得
\[
V_{\alpha_1} \bigcap V_{\alpha_2} \subset V_{\alpha_1}.
\]
从而
\[
\dim \left( V_{\alpha_1} \bigcap V_{\alpha_2} \right) < \dim \left( V_{\alpha_1} \right)
\]
此时要么
\[
\dim \left( V_{\alpha_1} \bigcap V_{\alpha_2} \right) > \dim \bigcap_{\alpha \in \Lambda} V_{\alpha}
\]
要么
\[
\dim \left( V_{\alpha_1} \bigcap V_{\alpha_2} \right) = \dim \bigcap_{\alpha \in \Lambda} V_{\alpha}
\]
因此我们可以继续讨论 \(\alpha_3 \in \Lambda \setminus \{\alpha_1, \alpha_2\}\)，依次下去.可以得到一个严格递减正整数列$\left\{ \mathrm{dim}\bigcap_{i=1}^n{V_{\alpha _i}} \right\} $,满足
\[
\mathrm{dim}\bigcap_{\alpha \in \Lambda}{V_{\alpha _i}}\leqslant \mathrm{dim}\bigcap_{i=1}^m{V_{\alpha _i}}\leqslant \mathrm{dim}V_{\alpha _1}，\mathrm{dim}\bigcap_{i=1}^m{V_{\alpha _i}}\in \mathbb{N},m\in \mathbb{N}.
\]
因此严格递减正整数列$\left\{ \mathrm{dim}\bigcap_{i=1}^n{V_{\alpha _i}} \right\}$只能有限项，因此这样的操作会经过有限次而停止.于是我们知道存在 \(\alpha_1, \alpha_2, \cdots, \alpha_n \in \Lambda\) 使得
\[
\bigcap_{i = 1}^n V_{\alpha_i} = \bigcap_{\alpha \in \Lambda} V_{\alpha}.
\]
\end{proof}

\begin{example}
设\( P \)为数域，\( \tau \)为\( P^{n \times n} \)上的线性变换，满足条件：对任意固定的\( A, B \in P^{n \times n} \)，\( \tau(AB) = \tau(A)\tau(B) \)或\( \tau(AB) = \tau(B)\tau(A) \)至少有一个成立，证明：要么对所有的\( A, B \in P^{n \times n} \)，\( \tau(AB) = \tau(A)\tau(B) \)，要么对所有的\( A, B \in P^{n \times n} \)，\( \tau(AB) = \tau(B)\tau(A) \).
\end{example}
\begin{proof}
对$\forall A\in P^{n\times n},$记
$$U_A\triangleq \left\{ B:\tau \left( AB \right) =\tau \left( A \right) \tau \left( B \right) \right\}, \quad U_{A}^{\prime}\triangleq \left\{ B:\tau \left( AB \right) =\tau \left( B \right) \tau \left( A \right) \right\},$$
则由条件知
$$P^{n\times n}=U_A\cup U_{A}^{\prime}.$$
从而$U_A,U_{A}^{\prime}$都不是非平凡子空间,否则与\hyperref[theorem:覆盖定理]{覆盖定理}矛盾!因此对$\forall A\in P^{n\times n},$要么$U_A=P^{n\times n},$要么$U_{A}^{\prime}=P^{n\times n}.$再记
$$W_1=\left\{ A:U_A=P^{n\times n} \right\}, \quad W_2=\left\{ A:U_{A}^{\prime}=P^{n\times n} \right\},$$
则
$$P^{n\times n}=W_1\cup W_2.$$
于是$W_1,W_2$都不是非平凡子空间,否则与\hyperref[theorem:覆盖定理]{覆盖定理}矛盾!故要么$W_1=P^{n\times n},$要么$W_2=P^{n\times n}.$即
\begin{align*}
\text{要么对}\forall A\in P^{n\times n},\text{都有}U_A=P^{n\times n};
\\
\text{要么对}\forall A\in P^{n\times n},\text{都有}U_{A}^{\prime}=P^{n\times n}.
\end{align*}
也即
\begin{align*}
\text{要么对}\forall A,B\in P^{n\times n},\text{都有}\tau \left( AB \right) =\tau \left( A \right) \tau \left( B \right) ;
\\
\text{要么对}\forall A,B\in P^{n\times n},\text{都有}\tau \left( AB \right) =\tau \left( B \right) \tau \left( A \right) .
\end{align*}
\end{proof}






\end{document}